\cotefile{29}
\scanpage{23}
% archive sheet number: 22
base $T$ ''. When $T$ is regular, the notions 1\textsuperscript{\uline{o}} and 2\textsuperscript{\uline{o}} coincide. The specialization theorem for the tame fundamental group can be proved, by the technique of SGA 1, when one uses the notion of tamely ramified covering in the broad sense, relative to the\unsure{divisors|divisor} $D_{\bar{t}}$ \inserted{in the $S_{\bar{t}}$} ($D$ having normal crossings relative to $T$, $S$ proper and smooth over $T$).

\uline{Remarks} Obviously, for every $S'$ object of $\underline{\mathrm{Et}}_S$, one defines in an obvious way $\operatorname{Rev}_{\deleted{\underline{R}}}^{\inserted{\underline{R}}}(S')$, by considering the restriction of $\underline{R}$ to $\underline{\mathrm{Et}}_{S'}\quad(\underline{\mathrm{Et}}_S)_{/S'}$. For variable $S'$, one thus obtains a fibered category $\underline{\mathrm{Rev}}^{\underline{R}}$ over $\underline{\mathrm{Et}}_S$, whose fibers $\underline{\mathrm{Rev}}_{S'}^{\underline{R}}$ are the $\operatorname{Rev}^{\underline{R}}(S')$, and a cartesian functor
\[
\underline{r}:\quad\underline{\mathrm{Rev}}^{\underline{R}}\qquad\underline{\mathrm{Rev}},
\tag*{\unsure{(*)|(\deleted{=}*)|(\deleted{*}=)}}
\]
where for short one denotes by $\underline{\mathrm{Rev}}$ the category \inserted{fibered over $\underline{\mathrm{Et}}_S$} of coverings of the different objects $S'$ of $\underline{\mathrm{Et}}_S$. (NB one assumes that for every $S'$, the geometric realization functor on $\underline{C}_{S'}$ takes values in the coverings of $S'$, and that the geometric realization functor on $\operatorname{Rev}^{\underline{R}}(S')$ is well defined). One must insist on the fact that for us, the ``ramification datum'' $\underline{R}$ is an object of an essentially auxiliary nature, intended to define the fibered category $\underline{\mathrm{Rev}}^{\underline{R}}$ (which is even a \uline{\unsure{s\deleted{t}a\deleted{a}tck|s\deleted{a}t\deleted{t}ack}}) and the cartesian functor $\underline{r}$ of (*); it is \unsure{(*)|(\deleted{=}*)|(\deleted{*}=)} that is the true object of study. One will therefore regard two ramification data $\underline{R}$, $\underline{R}'$ as essentially equivalent, if they define equivalent functors (*) (by an equivalence, of course, inducing the identity on $\underline{\mathrm{Rev}}$). One can convince oneself from examples that $\underline{R}$ and $\underline{R}'$ can have very different appearances. Thus, under the conditions of example 3), if $\underline{R}'$ is defined by a fibered subcategory \deleted{of} $\underline{C}'$ of $\underline{C}$
